\BourbakiExercises{习题}

\begin{enumerate}
\def\labelenumi{\arabic{enumi})}
\item
  设 A 是一个交换环，G 是一个群，H 是 G 的一个有限指标的子群。假设 A 的元素 \((G:H)1_{A}\) 可逆。设 M 是一个 A{[}G{]}-模，N 是 M 的一个 A{[}G{]}-子模。证明：若 N 在 M 中容许一个补的 A{[}H{]}-子模，则它在 M 中容许一个补的 A{[}G{]}-子模（仿照 VIII, 第 401 页定理 1 的证明）。由此推出，一个作为 A{[}H{]}-模为投射模的 A{[}G{]}-模是投射模。
\item
  设 A 是一个交换环，G 是一个群。
\end{enumerate}

\begin{enumerate}
\def\labelenumi{\alph{enumi})}
\item
  假设群 G 的代数 A{[}G{]} 是一个绝对平坦环（VIII, 第 171 页，习题 27）。证明环 A 是绝对平坦的。设 H 是 G 的一个容许有限生成子集的子群；证明左理想 \(I_{H}\)（A{[}G{]} 的左理想，VIII, 第 22 页，习题 25）在 A{[}G{]} 中容许一个补，并由此推出 H 是有限的（注意，在相反的情形，\(I_{H}\) 的零化子将化为 (0)）。
\item
  保留 a) 的假设。设 g 是 G 的一个阶为有限数 n 的元素。证明 \(n1_{A}\) 在 A 中可逆（设 x 是 A 中满足 \((1-e_{g})x(1-e_{g})=(1-e_{g})\) 的元素；构造 A 的元素 y，使得 \(1-(1-e_{g})x=y(1+e_{g}+\cdots+e_{g^{n-1}})\)）。由此推出，G 的每个有限子群的阶都在 A 中可逆。
\item
  假设 A 是一个绝对平坦环，并且 G 的每个有限子集都生成一个基数在 A 中可逆的有限子群。证明环 A{[}G{]} 是绝对平坦的（归结为 G 有限的情形，并利用习题 1）。
\end{enumerate}

\begin{enumerate}
\def\labelenumi{\arabic{enumi})}
\setcounter{enumi}{2}
\tightlist
\item
  设 A 是一个交换环，G 是一个群。
\end{enumerate}

\begin{enumerate}
\def\labelenumi{\alph{enumi})}
\item
  假设群 G 的代数 A{[}G{]} 是半单的。证明 A 是半单的，G 是有限的，并且它的阶在 A 中可逆（利用习题 2 以及 VIII, 第 22 页的习题 25）。
\item
  假设代数 A{[}G{]} 同构于一个有限域族的乘积。证明 G 的每个子群都是正规的（利用 a)，并注意 A{[}G{]} 中的每个幂等元都是中心的）。
\end{enumerate}

在下面的习题 4 至 26 中，G 是一个有限群，K 是一个代数闭体；我们假设 G 的阶在 K 中可逆。

\begin{enumerate}
\def\labelenumi{\arabic{enumi})}
\setcounter{enumi}{3}
\item
  设 \((\mathrm{M},\pi)\) 是 G 的一个忠实的 \(^{(2)}\) 有限维表示，\((\mathrm{N},\rho)\) 是一个不可约表示；假设特征标 \(\chi_{\pi}\) 恰好取 s 个不同的值。证明存在整数 n \textless{} s，使得 K{[}G{]}-模 N 同构于 \(M^{\otimes n}\) 的一个子模（计算 \(\langle\chi_{\rho},\chi_{\pi}^{n}\rangle\)，并注意等式 \(\chi_{\pi}(g)=\dim M\) 等价于 g=1）。
\item
  设 \((\mathrm{M},\pi)\) 是 G 的一个表示。
\end{enumerate}

\begin{enumerate}
\def\labelenumi{\alph{enumi})}
\item
  设 H 是 G 的一个子群，S 是 G 中（模 H 的）左陪集的一个代表系，N 是 M 的一个在 H 下稳定的线性子空间。G 在 M 中的表示同构于由 H 在 N 中的表示诱导的表示，当且仅当 \(M = \bigoplus_{s \in S} sN\)。
\item
  假设 M 是子空间族 \((\mathrm{M}_{i})_{i\in\mathrm{I}}\) 的直和，并且存在 G 在 I 上的一个传递作用，使得 \(\pi(g)(\mathrm{M}_{i})\subset\mathrm{M}_{gi}\)（对 \(g\in G\)、\(i\in I\)）。设 \(o\in I\)，\(G_{o}\) 是它在 G 中的稳定子。证明 \(\pi\) 同构于由 \(G_{o}\) 在 \(M_{o}\) 中的表示诱导的表示。若 \(\pi\) 不可约，则 G 在 I 上的作用是传递的这一条件自动满足。
\end{enumerate}

\begin{enumerate}
\def\labelenumi{\arabic{enumi})}
\setcounter{enumi}{5}
\tightlist
\item
  设 X 是一个 G 作用于其上的有限集。我们考虑表示 \(\pi\)（G 在 \(K^{X}\) 中的表示），它满足 \(\pi(g)e_{x}=e_{gx}\)（对 \(g\in G\)、\(x\in X\)）（与 X 相伴的“置换表示”）。
\end{enumerate}

\begin{enumerate}
\def\labelenumi{\alph{enumi})}
\item
  对 \(g \in G\)，证明 \(\chi_{\pi}(g)\) 等于 X 中被 g 固定的点的个数。
\item
  证明单位表示在 \(\pi\) 中的重数等于 G 在 X 中的轨道的个数。
\item
  从现在起假设 G 的作用是传递的；取定 X 的一点 x，并用 H 表示它的稳定子。证明 \(\pi\) 同构于由 H 的单位表示诱导的表示，并且 \(\langle\chi_{\pi},\chi_{\pi}\rangle\) 等于 H 在 X 中的轨道的个数。
\item
  设 V 是 \(K^{X}\) 中由诸坐标之和为零的向量组成的子空间；它在 G 下稳定，我们用 \(\pi_{X}\) 表示由 \(\pi\) 得出的 G 在 V 中的表示。证明 \(\pi_{X}\) 不可约当且仅当 G 的作用是双传递的（I, §5, 第 143 页，习题 14）。
\end{enumerate}

\begin{enumerate}
\def\labelenumi{\arabic{enumi})}
\setcounter{enumi}{6}
\item
  对 G 的每个有限维表示 V，用 \(\varphi(\mathrm{V})\) 表示 V 中被 G 固定的子空间的维数；我们把 \(\varphi\) 线性延拓为 \(\mathrm{R}_{\mathrm{K}}(\mathrm{G})\) 上的一个 Z-线性型。设 \((\mathrm{V}_{\lambda})_{\lambda\in\widehat{\mathrm{G}}}\) 是 G 的不可约表示族，\(\omega\) 是元素 \(\sum_{\lambda\in\widehat{\mathrm{G}}}\left[\mathrm{V}_{\lambda}\right]\left[\mathrm{V}_{\lambda}^{*}\right]\)（\(\mathrm{R}_{\mathrm{K}}(\mathrm{G})\) 的元素）。证明公式 \(\varphi(\omega^{n})=\sum_{\mathrm{C}}d(\mathrm{C})^{n-1}\) 对每个整数 \(n\geqslant0\) 成立（应用 VIII, 第 415 页的命题 8 以及特征标的第二正交关系）。
\item
  设 H 与 F 是 G 的子群，S 是 G 的一个子集，使得 G 是诸双陪集 FsH 的不交并。设 \((N,\sigma)\) 是 H 的一个表示。对 \(s\in G\)，令 \(H_{s}=F\cap sHs^{-1}\)，并用 \(\sigma_{s}\) 表示 \(H_{s}\) 在 N 中的由 \(\sigma_{s}(x)=\sigma(s^{-1}xs)\) 定义的表示。
\end{enumerate}

\begin{enumerate}
\def\labelenumi{\alph{enumi})}
\item
  证明表示 \(\operatorname{Res}_{\mathrm{F}}^{\mathrm{G}}(\operatorname{Ind}_{\mathrm{H}}^{\mathrm{G}}(\sigma))\) 同构于诸表示 \(\operatorname{Ind}_{\mathrm{H}_{s}}^{\mathrm{F}}(\sigma_{s})\)（\(s \in S\)）的直和。（把表示 \(\operatorname{Ind}_{\mathrm{H}}^{\mathrm{G}}(\sigma)\) 记作 \((\mathrm{M}, \pi)\)。空间 M 可与诸 \(\pi(x)\mathrm{N}\)（\(x \in G/H\)）的直和等同。对 \(s \in S\)，设 \(\mathrm{M}(s)\) 是 M 中由诸 \(\pi(x)\mathrm{N}\)（\(x \in FsH\)）生成的子空间。注意 \(\mathrm{M}(s)\) 是诸子空间 \(\pi(xs)\mathrm{N}\)（\(x \in F/H_{s}\)）的直和，并由习题 5, b) 推出：作为 K{[}F{]}-模，\(\mathrm{M}(s)\) 同构于 \(\operatorname{Ind}_{\mathrm{H}_{s}}^{\mathrm{F}}(\sigma_{s})\)。
\item
  取 F = H（于是我们有 \(H_{s} = H \cap sHs^{-1}\)）。证明表示 \(\operatorname{Ind}_{\mathrm{H}}^{\mathrm{G}}(\sigma)\) 不可约，当且仅当 \(\sigma\) 不可约，并且对每个 \(s \in G - H\)，表示 \(\sigma_{s}\) 与 \(\operatorname{Res}_{\mathrm{H}_{s}}^{\mathrm{H}}(\sigma)\) 的支集（VIII, 第 66 页）不相交（“麦基判别法”：应用 a) 与弗罗贝尼乌斯互反律）。
\item
  假设 H 是正规的；这时 \(\operatorname{Ind}_{\mathrm{H}}^{\mathrm{G}}(\sigma)\) 不可约当且仅当 \(\sigma\) 不可约，且不与其共轭 \(\sigma \circ \operatorname{int}(g)\)（\(g \in G - H\)）中的任何一个同构。
\end{enumerate}

\begin{enumerate}
\def\labelenumi{\arabic{enumi})}
\setcounter{enumi}{8}
\tightlist
\item
  设 \((\mathrm{M},\pi)\) 是 G 的一个不可约表示，H 是 G 的一个正规子群。a) 设 \(M=\bigoplus_{i\in I}M_{i}\) 是 K{[}H{]}-模 M 分解为同型子模直和的分解。证明 G 在 M 上的作用传递地置换诸子模 \(M_{i}\)。由此推出 \((\mathrm{M},\pi)\) 同构于由 G 的一个包含 H 的子群诱导的表示（应用习题 5, b)）。
\end{enumerate}

\begin{enumerate}
\def\labelenumi{\alph{enumi})}
\setcounter{enumi}{1}
\tightlist
\item
  设 \(\mathcal{S}\) 是 K{[}H{]}-模 M 的支集；存在整数 e，使得 M（作为 K{[}H{]}-模）同构于 \(\sum_{\lambda\in \mathcal{S}}\lambda^{e}\)。
\end{enumerate}

\begin{enumerate}
\def\labelenumi{\arabic{enumi})}
\setcounter{enumi}{9}
\tightlist
\item
  设 \((\mathrm{M},\pi)\) 是 G 的一个不可约表示，H 是 G 的一个正规子群。a) 假设 H 等于 G 的中心。证明 \(\pi\) 的次数整除 H 在 G 中的指标。（对 \(m \geqslant 1\)，设 \(H_{m}\) 是 \(H^{m}\) 中由元素 \((h_{1},\ldots,h_{m})\)（满足 \(h_{1}\ldots h_{m}=1\) 的元素）组成的子群。注意表示 \(\pi^{\times m}:G^{m}\to\operatorname{End}_{K}(M^{\otimes m})\) 定义了 \(G^{m}/H_{m}\) 的一个不可约表示，并对它应用命题 10 的推论 2（VIII, 第 420 页）。）
\end{enumerate}

\begin{enumerate}
\def\labelenumi{\alph{enumi})}
\setcounter{enumi}{1}
\tightlist
\item
  假设 H 是交换的。证明 \(\pi\) 的次数整除 (G : H)。（利用习题 9, a)，并对 G 的基数作归纳，归结到 \(\pi\) 在 H 上的限制是同型的情形。证明 \(\pi(\mathrm{H})\) 在 \(\pi(\mathrm{G})\) 中是中心的，并应用 a)。）
\end{enumerate}

\begin{enumerate}
\def\labelenumi{\arabic{enumi})}
\setcounter{enumi}{10}
\tightlist
\item
  设 A 与 H 是 G 的子群。假设 A 是交换的且是正规的，并且 G 是 H 与 A 的半直积。群 H 作用在 \(\widehat{\mathrm{A}} = \mathrm{Hom}(\mathrm{A}, \mathrm{K}^{*})\) 上；取定代表系 \((\alpha_{i})_{i \in \widehat{\mathrm{A}}/\mathrm{H}}\)（H 在 \(\widehat{A}\) 中的轨道的代表系）。对 \(i \in \widehat{A}/H\)，设 \(H_{i}\) 是 \(\alpha_{i}\) 在 H 中的稳定子，并令 \(G_{i} = AH_{i}\)。对每个不可约表示 \(\sigma\)（\(H_{i}\) 的不可约表示），定义表示 \(\rho_{i,\sigma}\)（\(G_{i}\) 的一个表示）：令 \(\rho_{i,\sigma}(ah) = \alpha_{i}(a)\sigma(h)\)（对 \(a \in A\) 与 \(h \in H\)）。证明诸表示 \(\operatorname{Ind}_{G_{i}}^{G}(\rho_{i,\sigma})\)（\(i \in \widehat{A}/H\)，\(\sigma \in \widehat{H}_{i}\)）都是不可约的且两两不同构，并且用这种方法我们得到 G 的全部不可约表示（应用习题 9）。
\end{enumerate}

¶ 12) 设 H 是 G 的一个子群；假设对每个 \(g \in G - H\)，我们有等式 \(H \cap gHg^{-1} = \{1\}\)。

\begin{enumerate}
\def\labelenumi{\alph{enumi})}
\item
  设 u 是 H 上的一个中心函数。证明函数 \(\operatorname{Ind}_{\mathrm{H}}^{\mathrm{G}}(u)\) 在 \(H - \{1\}\) 上与 u 一致。若 v 是 H 上另一个满足 \(v(1) = 0\) 的中心函数，则我们有 \(\langle \operatorname{Ind}_{\mathrm{H}}^{\mathrm{G}}(u), \operatorname{Ind}_{\mathrm{H}}^{\mathrm{G}}(v) \rangle_{\mathrm{G}} = \langle u, v \rangle_{\mathrm{H}}\)。
\item
  设 \(\lambda \in \widehat{\mathrm{H}}\)；用 \(\chi_{\lambda}^{\prime}\) 表示 H 上的函数 \(\chi_{\lambda} - d_{\lambda}\)，用 \(\theta_{\lambda}\) 表示 G 上的函数 \(d_{\lambda} + \mathrm{Ind}_{\mathrm{H}}^{\mathrm{G}}(\chi_{\lambda}^{\prime})\)。证明 \(\theta_{\lambda}\) 是 G 的一个次数为 \(d_{\lambda}\)、在 H 上的限制同构于 \(\lambda\) 的不可约表示的特征标（利用 a) 计算 \(\langle\theta_{\lambda},\theta_{\lambda}\rangle\)、\(\langle\theta_{\lambda},1\rangle\) 与 \(\langle\mathrm{Res}_{\mathrm{H}}^{\mathrm{G}}(\theta_{\lambda}),\chi_{\lambda}\rangle\)）。
\item
  令 \(\theta = \sum_{\lambda} d_{\lambda} \theta_{\lambda}\)。证明：\(\theta(h) = 0\)（若 \(h \in \mathrm{H} - \{1\}\)），且 \(\theta(g) = \theta(1) = \mathrm{Card}(\mathrm{H})\)（若 \(g\) 不属于 \(\mathrm{H}\) 的任何共轭）。
\item
  设 \(L'\) 是 G 中不属于 H 的任何共轭的元素所成的集合，并令 \(L = L' \cup \{1\}\)。证明 L 是 G 的一个正规子群（由 c) 推出 L 是一个以 \(\theta\) 为特征标的表示的核）。
\item
  证明 G 是 H 与 L 的半直积（计算 L 的阶）。f) 设 \(h \in H - \{1\}\)。证明自同构 \(\operatorname{int}(h)\) 在 L 上的限制除 1 外没有不动点。由此推出 H 的阶整除 \(\operatorname{Card}(L) - 1\)。
\end{enumerate}

\(g\) ) 证明：若子群 H 的阶为偶数，则 L 是交换的（应用 I, §6, 第 145 页的习题 23, \(d\) )）。

\begin{enumerate}
\def\labelenumi{\arabic{enumi})}
\setcounter{enumi}{12}
\tightlist
\item
  设 \(\mathcal{H}\) 是 \(G\) 的子群所成的一个集合。
\end{enumerate}

\begin{enumerate}
\def\labelenumi{\alph{enumi})}
\tightlist
\item
  证明下列性质等价：
\end{enumerate}

\begin{enumerate}
\def\labelenumi{(\roman{enumi})}
\item
  属于 \(\mathcal{H}\) 的诸子群的共轭之并等于 G。
\item
  \(G\) 的每个特征标都是属于 \(\mathcal{H}\) 的子群的特征标所诱导的特征标的有理系数线性组合。
\end{enumerate}

（由 VIII, 第 412 页的命题 6 推出，性质 (ii) 等价于限制同态 \(\mathcal{Z}_{\mathrm{K}}(\mathrm{G})\to \bigoplus_{\mathrm{H}\in \mathcal{H}}\mathcal{Z}_{\mathrm{K}}(\mathrm{H}).\) 的单性。）

\begin{enumerate}
\def\labelenumi{\alph{enumi})}
\setcounter{enumi}{1}
\item
  设 \(\mathcal{C}\) 是 G 的循环子群所成的集合；它显然具有性质 (i)。对 G 的每个循环子群 C，用 \(\varphi_{\mathrm{C}}\) 表示 C 上的一个函数，它在 C 的生成元上取值 \(\operatorname{Card}(\mathrm{C})\)，在其他元素上取值 0。证明等式 \(\operatorname{Card}(\mathrm{G}) = \sum_{\mathrm{C} \in \mathcal{C}} \operatorname{Ind}_{\mathrm{C}}^{\mathrm{G}}(\varphi_{\mathrm{C}})\)。
\item
  设 \(C \in \mathcal{C}\)。证明函数 \(\varphi_{C}\) 属于子群 \(\mathrm{R}_{\mathrm{K}}(\mathrm{G})\)（\(\mathcal{Z}_{\mathrm{K}}(\mathrm{G})\) 的由诸特征标生成的子群）（对 C 的阶用归纳法，并利用 b)）。
\end{enumerate}

\(d\) ) 设 \(\chi\) 是 G 的一个不可约表示的特征标。证明等式 \(\chi = \mathrm{Card}(\mathrm{G})^{-1}\sum_{\mathrm{C}\in \mathcal{C}}\mathrm{Ind}_{\mathrm{C}}^{\mathrm{G}}(\varphi_{\mathrm{C}}\mathrm{Res}_{\mathrm{C}}^{\mathrm{G}}(\chi))\)，它给出了 \(a\) ) 中性质 (ii) 的一个显式公式。

\begin{enumerate}
\def\labelenumi{\arabic{enumi})}
\setcounter{enumi}{13}
\tightlist
\item
  若 G 的一个表示是它的子群的 1 次表示所诱导的表示的直和，则称它是单项表示。
\end{enumerate}

\begin{enumerate}
\def\labelenumi{\alph{enumi})}
\item
  假设 G 的真子群的表示都是单项的，并且 G 包含一个非中心的交换正规子群 H。证明 G 的每个忠实的不可约表示都是单项的（注意这样一个表示在 H 上的限制不是同型的，并应用习题 9, a)）。
\item
  由此推出，超可解群（I, §6, 第 146 页，习题 26），特别是幂零群的所有表示都是单项的。
\item
  取 G 为群 \(\mathbf{SL}(2,\mathbf{F}_{3})\)。它的中心 Z 的阶为 2，并且群 G/Z 同构于交错群 \(\mathfrak{A}_{4}\)（II, §10, 第 422 页，习题 14, g)）；特别地，G 是可解的。证明它的导群的指标为 3（参见 I, §5, 第 142 页，
\end{enumerate}

习题 10）。试由 VIII, 第 407 页的公式 (8) 推出 G 有一个次数为 2 的不可约表示，并证明该表示不是单项的。

¶ 15) 假设 G 的所有表示都是单项的；我们来证明 G 是可解的。对 G 的阶作归纳法，我们可以假设 G 的每个异于 G 的商群都是可解的。

\begin{enumerate}
\def\labelenumi{\alph{enumi})}
\tightlist
\item
  若 G 包含两个不同的非平凡极小正规子群 \(H_{1}\) 与 \(H_{2}\)，则 G 是可解的（考虑从 G 到 \(G/H_{1} \times G/H_{2}\) 的典范同态）。
\end{enumerate}

从现在起假设 G 有唯一的非平凡极小正规子群 H。

\begin{enumerate}
\def\labelenumi{\alph{enumi})}
\setcounter{enumi}{1}
\tightlist
\item
  证明：若 G 的一个表示在 H 上的限制是非平凡的，则它是忠实的。c) 设 \(\pi\) 是一个极小次数的忠实表示；存在 G 的子群 \(H_{i}\) 与表示 \(\sigma_{i}\)（\(H_{i}\) 的次数为 1 的表示），使得 \(\pi\) 同构于 \(\bigoplus\operatorname{Ind}_{\operatorname{H}_{i}}^{\operatorname{G}}(\sigma_{i})\)。令 \(\pi_{0}=\bigoplus\operatorname{Ind}_{\operatorname{H}_{i}}^{\operatorname{G}}(\varepsilon_{\operatorname{H}_{i}})\)，其中 \(\varepsilon_{H_{i}}\) 表示单位表示。证明 \(\pi_{0}\) 的核是交换的且非平凡（注意 \(\pi_{0}\) 是不可约的）；由此推出 G 是可解的。
\end{enumerate}

\begin{enumerate}
\def\labelenumi{\arabic{enumi})}
\setcounter{enumi}{15}
\tightlist
\item
  假设 K 的特征为零。我们说 K 的元素 x 在 Z 上是整的，如果它属于 K 的一个作为有限生成 Z-模的子环，或者等价地，如果 K 的子环 \(Z[x]\) 是一个有限生成的 Z-模 \(^{*}\)（参见 Comm. Alg., V, §1, No.~1, 第 303 页）。
\end{enumerate}

\begin{enumerate}
\def\labelenumi{\alph{enumi})}
\item
  证明 K 中在 Z 上是整的元素所成的集合是 K 的一个子环（对在 Z 上都是整的 x, y，考虑子环 \(Z[x, y]\)）。
\item
  元素 \(x\)（\(K\) 的元素）在 \(\mathbf{Z}\) 上是整的，当且仅当它是 \(\mathbf{Z}[X]\) 中某个首一多项式的根（当 \(x\) 是整的时候，仿照 VIII, 第 81 页引理 1 的证明，证明它满足形如 \(\det(xI_m - A) = 0\) 的方程，其中 \(A \in \mathbf{M}_m(\mathbf{Z})\)）。
\item
  若 x 在 Z 上是整的，则它的在 Q 上的共轭 \(x_{1},\ldots,x_{m}\) 也是整的；元素 \(x_{1}\cdots x_{m}\) 属于 Z（利用 VIII, 第 416 页的引理）。
\end{enumerate}

\begin{enumerate}
\def\labelenumi{\arabic{enumi})}
\setcounter{enumi}{16}
\tightlist
\item
  设 \(\mathrm{C}\) 是 \(\mathrm{G}\) 的一个共轭类，\(\pi\) 是 \(\mathrm{G}\) 的一个次数为 \(d\) 的不可约表示。
\end{enumerate}

\begin{enumerate}
\def\labelenumi{\alph{enumi})}
\item
  证明 K 的元素 \(d^{-1}\operatorname{Card}(\mathrm{C})\chi(\mathrm{C})\) 在 Z 上是整的（按 VIII, 第 415 页命题 9 的证明的方式推理）。
\item
  假设 d 与 \(\operatorname{Card}(\mathbf{C})\) 是互素的整数。证明元素 \(d^{-1}\chi(\mathbf{C})\) 在 Z 上是整的。
\end{enumerate}

* \(c\) ) 此外假设 \(\mathrm{K} = \mathbf{C}\)。证明 \(d^{-1}\chi (\mathrm{C})\) 在 \(\mathbf{Q}\) 上的所有共轭的绝对值都 \(\leqslant 1\)。由习题 16, \(c\) ) 推出 \(|\chi (\mathrm{C})|\) 或者为零或者等于 \(d\)。由此得出结论：若 \(\chi (\mathrm{C})\) 非零，则 \(\pi (x)\)（对每个 \(x\in \mathrm{C}\)）都是一个位似。特别地，若 \(\pi\) 是忠实的且 C 不化为单个元素，则我们有 \(\chi (\mathrm{C}) = 0.\ast\)

\begin{enumerate}
\def\labelenumi{\arabic{enumi})}
\setcounter{enumi}{17}
\tightlist
\item
  假设群 G 是单群但不是循环群。
\end{enumerate}

\begin{enumerate}
\def\labelenumi{\alph{enumi})}
\item
  证明不存在基数是素数幂且异于 \(\{1\}\) 的共轭类。（设 C 是一个基数为 \(p^{r}\)（\(r \geqslant 1\)）的共轭类。由习题 17, b) 推出，\(p^{-1}\chi_{\lambda}(1)\chi_{\lambda}(C)\)（当 \(\lambda\) 异于单位表示时）在 Z 上是整的，并用特征标的第二正交关系导出矛盾。）
\item
  证明 G 的阶至少被三个不同的素数整除（对 G 的一个西罗子群中的中心元素 \(x \neq 1\) 的共轭类应用 a)）。
\item
  证明阶至多被两个不同素数整除的有限群是可解的（“伯恩赛德定理”：对群的基数用归纳法，并利用 b)）。
\end{enumerate}

*19) 设 \(\pi\) 是 G 的一个次数 \(>1\) 的不可约复表示。证明存在 G 的一个元素 \(x\)，使得 \(\chi_{\pi}(x)=0\)。（设 \(\{\lambda_1,\ldots,\lambda_s\}\) 是 \(\pi\) 在 \(\widehat{G}\) 中的轨道（群 \(\operatorname{Gal}(\mathbf{C}/\mathbf{Q})\) 作用下的轨道），\(\chi_1,\ldots,\chi_s\) 是相应的特征标。设 \(g\in G\)；若 \(\chi(g)\neq 0\)，则仿照习题 17, \(c\) ) 的方式证明 \(|\chi_1(g)\ldots\chi_s(g)|\geqslant 1\)，从而 \(\sum_i|\chi_i(g)|^2\geqslant s\)（FRV, III, §1, No.~1, 第 93 页，命题 2）。若这对每个 \(g\in G\) 都成立，则由特征标的正交关系推出我们取等号，而取 \(g=1\) 就导出矛盾。）*

\begin{enumerate}
\def\labelenumi{\arabic{enumi})}
\setcounter{enumi}{19}
\tightlist
\item
  设 G 是一个有限群，H 是 G 的一个子群，\(\sigma\) 是 H 的一个有限维表示，\(\rho\) 是由 \(\sigma\) 诱导的 G 的表示。设 C 是 G 的一个共轭类；集合 \(C \cap H\) 是 H 的共轭类的一个有限族 \((D_{i})_{i \in I}\) 的不交并。
\end{enumerate}

\begin{enumerate}
\def\labelenumi{\alph{enumi})}
\item
  证明公式 \(\chi_{\rho}(\mathrm{C}) = \frac{(\mathrm{G:H})}{\mathrm{Card}(\mathrm{C})}\sum_{i\in \mathrm{I}}\mathrm{Card}(\mathrm{D}_i)\chi_\sigma (\mathrm{D}_i)\)。
\item
  若 \(\sigma\) 是平凡表示，则有 \(\chi_{\rho}(\mathrm{C})=(\mathrm{G}:\mathrm{H})\operatorname{Card}(\mathrm{C})^{-1}\operatorname{Card}(\mathrm{C}\cap\mathrm{H})\)。
\end{enumerate}

¶ 21) 设 \(n \in \mathbf{N}\)。用 \(\mathcal{D}_{n}\) 表示 \(\mathbf{N}^{n}\) 中由元素 \((p_{1},\ldots,p_{n})\)（满足 \(p_{1} \geqslant \cdots \geqslant p_{n}\) 的元素）组成的子集；赋予它字典序。对 \(\mathbf{p} \in \mathcal{D}_{n}\)，用 \(\mathbf{p}'\) 表示 \(\mathcal{D}_{n}\) 中由 \(p_{i}' = p_{i} + n - i\) 定义的元素。设 \(\mathbf{X} = (\mathrm{X}_{1},\ldots,\mathrm{X}_{n})\) 是一个变量序列；令 \(\Delta(\mathbf{X}) = \prod_{i<j}(\mathrm{X}_{i} - \mathrm{X}_{j})\)。

\begin{enumerate}
\def\labelenumi{\alph{enumi})}
\item
  设 \(\mathrm{P} \in \mathbf{Z}[\mathbf{X}]\) 是一个反对称多项式，即满足关系 \(\mathrm{P}(\mathrm{X}_{\sigma(1)}, \ldots, \mathrm{X}_{\sigma(n)}) = \varepsilon(\sigma)\mathrm{P}(\mathrm{X}_1, \ldots, \mathrm{X}_n)\)（对每个元素 \(\sigma \in \mathfrak{S}_n\)）的多项式。证明存在对称多项式 \(\mathrm{Q} \in \mathbf{Z}[\mathbf{X}]\)，使得 \(\mathrm{P} = \Delta\mathrm{Q}\)。
\item
  对每个 \(\mathbf{p} \in \mathcal{D}_n\)，令 \(S_{\mathbf{p}}(\mathbf{X}) = \Delta(\mathbf{X})^{-1} \det(X_j^{p_i'})_{1 \leqslant i,j \leqslant n}\)。证明 \(S_{\mathbf{p}}\) 是一个次数为 \(\sum p_i\)、具有整系数的齐次对称多项式。
\item
  若 \(\mathbf{q} = (q_i)_{1\leqslant i\leqslant n}\) 是 \(\mathcal{D}_n\) 的一个元素，而 \(P\) 是 \(\mathbf{Q}[\mathbf{X}]\) 的一个元素，则用 \([P]_{\mathbf{q}}\) 表示单项式 \(\mathbf{x}^{\mathbf{q}}\) 在 \(P\) 中的系数。令 \(K_{\mathbf{pq}} = [S_{\mathbf{p}}]_{\mathbf{q}}\)。证明我们有 \(K_{\mathbf{pp}} = 1\) 与 \(K_{\mathbf{pq}} = 0\)（对 \(\mathbf{q} > \mathbf{p}\)）（按字典序展开等式 \(\det(X_j^{p_i'}) = S_{\mathbf{p}}(\mathbf{X})\Delta(\mathbf{X})\)）。由此推出，诸多项式 \(S_{p}\)（\(p \in \mathcal{D}_{n}\)）构成 Z{[}X{]} 中由对称多项式组成的 Z-子模的一组基。
\item
  设 \(\mathbf{Y} = (\mathrm{Y}_1, \ldots, \mathrm{Y}_n)\) 是一个变量序列。证明形式幂级数的等式
\end{enumerate}

\[
\prod_ {1 \leqslant i, j \leqslant n} (1 - \mathrm{X} _ {i} \mathrm{Y} _ {j}) ^ {- 1} = \sum_ {\mathbf {p} \in \mathcal {D} _ {n}} \mathrm{S} _ {\mathbf {p}} (\mathbf {X})   \mathrm{S} _ {\mathbf {p}} (\mathbf {Y})
\]

（利用等式 \(\det\big((1 - X_iY_j)^{-1}\big)_{1\leqslant i,j\leqslant n} = \Delta(\mathbf{X})\Delta(\mathbf{Y})\prod_{1\leqslant i,j\leqslant n}(1 - X_iY_j)^{-1}\)（参见 III, §8, 第 639 页，习题 5），把 \((1 - T)^{-1}\) 作为形式幂级数展开以计算左边）。

\begin{enumerate}
\def\labelenumi{\alph{enumi})}
\setcounter{enumi}{4}
\tightlist
\item
  设 \(\mathrm{P} \in \mathbf{Q}[\mathbf{X}]\) 是一个对称多项式。对每个 \(\mathbf{p} \in \mathcal{D}_n\)，令 \(\omega_{\mathbf{p}}(\mathrm{P}) = [\mathrm{P}\Delta]_{\mathbf{p}'}\)。证明等式 \(\omega_{\mathbf{p}}(\mathrm{S}_{\mathbf{q}}) = \delta_{\mathbf{pq}}\)。由此推出，对每个 \(\mathbf{q} \in \mathcal{D}_n\)，我们有 \([\mathrm{P}]_{\mathbf{q}} = \sum_{\mathbf{p} \in \mathcal{D}_n} \mathrm{K}_{\mathbf{pq}} \omega_{\mathbf{p}}(\mathrm{P})\)。
\end{enumerate}

\(f)\) 对 \(\boldsymbol{\alpha}\in\mathbf{N}^{n}\)，令 \(f^{\boldsymbol{\alpha}}(\mathbf{X})=\prod_{j=1}^{n}(\sum_{i}\mathrm{X}_{i}^{j})^{\alpha_{j}}\) 与 \(c_{\boldsymbol{\alpha}}=\prod_{j=1}^{n}\frac{1}{j^{\alpha_{j}}\alpha_{j}!}\)。证明我们有 \(\sum_{\boldsymbol{\alpha}\in\mathbf{N}^{n}}c_{\boldsymbol{\alpha}}\omega_{\mathbf{p}}(f^{\boldsymbol{\alpha}})\omega_{\mathbf{q}}(f^{\boldsymbol{\alpha}})=\delta_{\mathbf{pq}}\)（对 \(\mathbf{p},\mathbf{q}\)（\(\mathcal{D}_{n}\) 中的元素））（注意我们有

\[
\prod_ {1 \leqslant i, j \leqslant n} (1 - \mathrm{X} _ {i} \mathrm{Y} _ {j}) ^ {- 1} = \prod_ {k = 0} ^ {\infty} \exp \bigl (\frac {1}{k} (\sum_ {i} \mathrm{X} _ {i} ^ {k}) (\sum_ {j} \mathrm{Y} _ {j} ^ {k}) \bigr) = \sum_ {\boldsymbol {\alpha} \in \mathbf {N} ^ {n}} c _ {\boldsymbol {\alpha}} f ^ {\boldsymbol {\alpha}} (\mathbf {X}) f ^ {\boldsymbol {\alpha}} (\mathbf {Y}) \bigr).
\]

¶ 22) 设 n 是一个自然数。用 \(\mathcal{S}_{n}\) 表示和为 n 的严格正整数的递减有限序列所成的集合。设 \(\mathbf{p} = (p_{i})_{1 \leqslant i \leqslant d}\) 是 \(\mathcal{S}_{n}\) 的一个元素。对 \(1 \leqslant i \leqslant d\)，用 \(P_{i}\) 表示满足 \(p_{1} + \cdots + p_{i-1} < r \leqslant p_{1} + \cdots + p_{i}\) 的整数 r 所成的集合；对 \(1 \leqslant j \leqslant p_{1}\)，用 \(P_{j}'\) 表示形如 \(p_{1} + \cdots + p_{i-1} + j\)（其中 \(1 \leqslant i \leqslant d\) 且 \(p_{i} \geqslant j\)）的整数所成的集合。子集 \(P_{i}\) 一方面、子集 \(P_{j}'\) 另一方面，构成 N 中区间 {[}1, n{]} 的一个分拆。

（我们常用一个图来表示这一情形，该图称为杨图：在下面的例子中，我们有 n = 9 与 \(\mathbf{p} = (4, 2, 2, 1)\)。集合 \(P_{i}\) 由第 i 行的元素组成，而 \(P_{j}^{\prime}\) 由第 j 列的元素组成。）

1

2

3

4

5

6

7

8

9

设 \(\mathcal{P}\)（相应地，\(\mathcal{P}'\)）是 \(\mathfrak{S}_n\) 中由这样的置换 \(\sigma\) 组成的子群：\(\sigma(\mathrm{P}_i) \subset \mathrm{P}_i\)（对每个 \(i\)）且（相应地）\(\sigma(\mathrm{P}_j') \subset \mathrm{P}_j'\)（对每个 \(j\)）。

\begin{enumerate}
\def\labelenumi{\alph{enumi})}
\item
  证明我们有 \(\mathcal{P} \cap \mathcal{P}' = \{1\}\)。
\item
  设 \(\sigma \in \mathfrak{S}_n\)。证明：若 \(\mathrm{Card}(\mathrm{P}_i \cap \sigma(\mathrm{P}_j')) \leqslant 1\)（对每个 \(i\) 与每个 \(j\)），则 \(\sigma\) 属于 \(\mathcal{PP}'\)。由此推出，若 \(\sigma \notin \mathcal{PP}'\)，则子群 \(\mathcal{P} \cap \sigma \mathcal{P}'\sigma^{-1}\) 包含一个对换。
\item
  在群代数 \(\mathbf{Z}[\mathfrak{S}_n]\) 中，令
\end{enumerate}

\[
a _ {\mathbf {p}} = \sum_ {\sigma \in \mathcal {P}} e _ {\sigma}, b _ {\mathbf {p}} = \sum_ {\tau \in \mathcal {P} ^ {\prime}} \varepsilon (\tau) e _ {\tau}, c _ {\mathbf {p}} = a _ {\mathbf {p}} b _ {\mathbf {p}}.
\]

证明 \(c_{p}\) 满足 \(e_{\sigma}c_{p}e_{\sigma'}=\varepsilon(\sigma')c_{p}\)（对每个 \(\sigma\in \mathcal{P}\) 与每个 \(\sigma'\in \mathcal{P}'\)），并且具有这一性质的每个元素都属于 \(Zc_{p}\)（把这样一个元素写成 \(\sum n_{\sigma}e_{\sigma}\)，并由 b) 推出 \(n_{\sigma}\) 为零（当 \(\sigma\notin \mathcal{P}\mathcal{P}'\)））。由此推出，对 \(Z[\mathfrak{S}_n]\) 的每个元素 x，元素 \(c_{p}xc_{p}\) 属于 \(Zc_{p}\)。

\(d\) ) 用 A 表示 \(\mathbf{Q}\)-代数 \(\mathbf{Q}[\mathfrak{S}_n]\)。证明 A 的理想 \(V_{\mathbf{p}} = Ac_{\mathbf{p}}\) 是一个绝对单的 A-模。（注意我们有 \(c_{\mathbf{p}}V_{\mathbf{p}} \subset \mathbf{Q}c_{\mathbf{p}}\)，由 \(c\)。）证明：若 \(\mathfrak{a}\) 是一个真包含于 \(V_{\mathbf{p}}\) 之中的左理想，则我们有 \(c_{\mathbf{p}}\mathfrak{a} = 0\)，从而 \(\mathfrak{a}^2 = 0\)，最终 \(\mathfrak{a} = 0\)。）

\begin{enumerate}
\def\labelenumi{\alph{enumi})}
\setcounter{enumi}{4}
\tightlist
\item
  证明我们有 \(c_{\mathbf{p}}^{2} = (n!/[V_{\mathbf{p}} : \mathbf{Q}]) c_{\mathbf{p}}\)（计算 A 的自同态 \(x \mapsto xc_{p}\) 的迹）。
\end{enumerate}

¶ 23) 沿用上一习题的记号。设 \(\mathbf{p} = (p_i)\) 与 \(\mathbf{q} = (q_j)\) 是 \(S_n\) 的两个元素。

\begin{enumerate}
\def\labelenumi{\alph{enumi})}
\item
  假设按字典序有 \(\mathbf{p} > \mathbf{q}\)。证明我们有 \(a_{\mathbf{p}}xb_{\mathbf{q}} = 0\)（对每个 \(x \in \mathrm{A}\)）。（归结为 \(x = 1\) 的情形。注意若 \((\mathrm{P}_i), (\mathrm{P}_j')\) 表示与 \(\mathbf{p}\) 相伴的分拆，而 \((\mathrm{Q}_k), (\mathrm{Q}_l')\) 表示与 \(\mathbf{q}\) 相伴的分拆，则记 \(s\) 为最小的整数 \(n\)，使得 \(p_n \neq q_n\)，我们有 \(\operatorname{Card}(\mathrm{P}_s \cap \mathrm{Q}_1') \geqslant 2\)；构造一个对换 \(\tau\)，使得 \(a_{\mathbf{p}}\tau = a_{\mathbf{p}}\) 且 \(\tau b_{\mathbf{q}} = -b_{\mathbf{q}}\)。）
\item
  由此推出，若 \(p \neq q\)，则 A-模 \(V_{p}\) 与 \(V_{q}\) 不同构（应用 a)，并利用把 \(e_{\sigma}\) 映为 \(e_{\sigma^{-1}}\) 的 A 的反自同构）。
\item
  证明把序列 \(\mathbf{p}\) 映为表示 \((\mathrm{V}_{\mathbf{p}})_{(\mathrm{K})}\) 的映射定义了一个从 \(\mathcal{S}_n\) 到集合 \(\mathcal{S}_{\mathrm{K}}(\mathfrak{S}_n)\)（单 \(\mathrm{K}[\mathfrak{S}_n]\)-模的类所成的集合）的双射。d) 对 \(\mathbf{p} \in \mathcal{S}_n\)，用 \(\overline{\mathbf{p}}\) 表示由 \(\overline{p}_j = \mathrm{Card}(\mathrm{P}_j')\) 定义的序列（习题 22）。证明映射 \(\mathbf{p} \mapsto \overline{\mathbf{p}}\) 是 \(\mathcal{S}_n\) 的一个对合；A-模 \(\mathrm{V}_{\overline{\mathbf{p}}}\) 同构于 \(\mathrm{V}_{\mathbf{p}} \otimes_{\mathbf{Q}} \mathbf{Q}_{\varepsilon}\)，其中 \(\mathbf{Q}_{\varepsilon}\) 表示 \(\mathbf{Q}\) 上与符号相伴的一维表示。
\item
  描述与序列 \((1, \ldots, 1)\) 和 \((n)\) 相对应的表示，以及与序列 \((n - 1, 1)\) 相对应的表示。
\end{enumerate}

¶ 24) 沿用习题 21 至 23 的记号。设 \(\mathbf{p} = (p_i)_{1 \leqslant i \leqslant d} \in \mathcal{S}_n\)；用 \(\mathbf{p}\) 表示元素 \((p_1, \ldots, p_d, 0, \ldots, 0)\)（\(\mathcal{D}_n\) 的元素）。若 \(\boldsymbol{\alpha} = (\alpha_1, \ldots, \alpha_n)\) 是 \(\mathbf{N}^n\) 的满足 \(\sum_i i\alpha_i = n\) 的一个元素，则用 \(C_{\boldsymbol{\alpha}}\) 表示 \(\mathfrak{S}_n\) 的由如下元素组成的共轭类：这些元素是一族支集两两不交的轮换的乘积，其中 \(\alpha_2\) 个轮换的阶为 2，\(\alpha_3\) 个轮换的阶为 3，等等。

\begin{enumerate}
\def\labelenumi{\alph{enumi})}
\item
  证明表示 \(\rho_{\mathbf{p}}\)（\(\mathfrak{S}_n\) 在 A-模 \(\mathrm{Aa}_{\mathbf{p}}\) 中的表示）是由 \(\mathcal{P}\) 的平凡表示诱导的表示。利用习题 20 推出公式 \(\chi_{\rho_{\mathbf{p}}}(\mathrm{C}_{\alpha}) = [f^{\alpha}]_{\mathbf{p}}\)。
\item
  对任意 \(\mathbf{q} \in \mathcal{S}_n\)，用 \(\lambda_{\mathbf{q}}\) 表示 \(\mathfrak{S}_n\) 在 A-模 \(V_{\mathbf{q}}\) 中的表示。证明存在正整数 \(n_{\mathbf{pq}}\)，使得我们有 \(\chi_{\rho_{\mathbf{p}}} = \sum_{\mathbf{q} \in \mathcal{S}_n} n_{\mathbf{pq}} \chi_{\lambda_{\mathbf{q}}}\)（其中 \(n_{\mathbf{pp}} > 0\)）。
\item
  设 \(\eta_{\mathbf{p}}\) 是 \(\mathfrak{S}_n\) 上的取值 \(\omega_{\mathbf{p}}(f^{\alpha})\)（在类 \(\mathrm{C}_{\alpha}\) 的元素上）的中心函数。证明 \(\eta_{\mathbf{p}}\) 是诸特征标 \(\chi_{\lambda_{\mathbf{q}}}\)（\(\mathbf{q} \in \mathcal{S}_n\)）的具有整系数的线性组合（利用习题 21, c) 与 e)）。由此推出，\(\eta_{\mathbf{p}}\) 相差一个符号等于某个特征标 \(\chi_{\lambda_{\mathbf{q}}}\)。（利用习题 21, f) 计算 \(\langle \eta_{\mathbf{p}}, \eta_{\mathbf{p}} \rangle\)。）d) 推出等式 \(\chi_{\lambda_{\mathbf{p}}}(\mathrm{C}_{\alpha}) = [\Delta f^{\alpha}]_{\mathbf{p}'}\)（“弗罗贝尼乌斯公式”）。
\item
  证明不可约表示 \(\lambda_{p}\) 在表示 \(\rho_{q}\) 中的重数等于 \(K_{pq}\)。
\end{enumerate}

\(f\) ) 证明 \(\lambda_{\mathbf{p}}\) 的次数为 \(\frac{n!}{\mathbf{p}^{\prime}!}\Delta (\mathbf{p}^{\prime})\)（展开多项式 \(\Delta (\mathbf{X})(\sum \mathrm{X}_i)^n\)）。

¶ 25) 考虑表示 \(\pi_{\mathbf{n}}\)（\(\mathfrak{S}_n\) 的与 \(\mathfrak{S}_n\) 在集合 \(\mathbf{n} = [1, n]\) 上的作用相伴的表示）（习题 6）。设 \(k\) 是一个正整数，\(\chi_k\) 是表示 \(\wedge^k \pi_{\mathbf{n}}\) 的特征标。

\begin{enumerate}
\def\labelenumi{\alph{enumi})}
\item
  设 \(\sigma \in \mathfrak{S}_n\)。证明公式 \(\chi_k(\sigma) = \sum \varepsilon (\sigma |S)\)，其中求和取遍诸子集 \(S\)（\(\mathbf{n}\) 的具有 \(k\) 个元素且满足 \(\sigma (S) = S\) 的子集）所成的集合。
\item
  证明 \(\wedge^{k}\pi_{n}\) 是 \(\mathfrak{S}_n\) 的与元素 \((n-k,1,\ldots,1)\)（\(\mathcal{S}_n\) 的元素）相伴的一个不可约表示（利用弗罗贝尼乌斯公式计算这个表示的特征标，并注意和式
\end{enumerate}

\[
\Delta (\mathbf {X}) f ^ {\boldsymbol {\alpha}} (\mathbf {X}) = \sum_ {\sigma \in \mathfrak {S} _ {n}} \varepsilon (\sigma) X _ {1} ^ {\sigma (n) - 1} \dots X _ {n} ^ {\sigma (1) - 1} f ^ {\boldsymbol {\alpha}} (\mathbf {X})
\]

中 \(X^{p'}\) 的系数非零的那些项对应于这样的置换 \(\sigma\)：对每个 i 满足 \(\sigma(i) \leqslant i + 1\)，且满足 \(\sigma(i) = i\)（对 \(1 \leqslant i < n - k\)）。

¶ 26) 设 q 是一个素数的幂，F 是有 q 个元素的有限域；取 G 为群 \(\mathbf{GL}(2, \mathbf{F})\)。

\begin{enumerate}
\def\labelenumi{\alph{enumi})}
\item
  描述 G 的共轭类（区分四种类型的类，它们分别含有 1、\(q^{2}-1\)、\(q^{2}+q\) 与 \(q^{2}-q\) 个元素）。
\item
  群 G 作用在 F 上的射影直线 P 上；考虑 G 的表示 \(\pi_{P}\)（习题 6）。对任意群同态 \(\lambda : F^{*} \to K^{*}\)，用 \(\delta_{\lambda}\) 表示与同态 \(\lambda \circ det\)（从 G 到 \(K^{*}\) 的同态）相伴的次数为 1 的表示，并令 \(\pi_{\lambda} = \pi_{P} \otimes \delta_{\lambda}\)。证明 \(\pi_{\lambda}\) 不可约，并计算它的特征标。
\item
  设 B 是由上三角矩阵组成的 G 的子群。对 \(\lambda,\mu\)（\(\mathrm{Hom}(\mathrm{F}^{*},\mathrm{K}^{*})\) 中的元素），用 \(\beta_{\lambda,\mu}\) 表示由 \(\beta_{\lambda,\mu}((a_{ij}))=\lambda(a_{11})\mu(a_{22})\) 定义的 B 的次数为 1 的表示，并令 \(\pi_{\lambda,\mu}=\mathrm{Ind}_{\mathrm{B}}^{\mathrm{G}}(\beta_{\lambda,\mu})\)。证明 \(\pi_{\lambda,\mu}\) 当 \(\lambda\neq\mu\) 时不可约，而 \(\pi_{\lambda,\lambda}\) 同构于 \(\pi_{P}\oplus\delta_{\lambda}\)（应用习题 8）。表示 \(\pi_{\lambda,\mu}\) 与 \(\pi_{\lambda',\mu'}\) 同构，当且仅当子集 \(\{\lambda,\mu\}\) 与 \(\{\lambda',\mu'\}\)（\(\mathrm{Hom}(\mathrm{F}^{*},\mathrm{K}^{*})\) 的子集）相等。
\end{enumerate}

\(d\) ) 设 \(\mathrm{F}^{\prime}\) 是 F 的一个次数为 2 的扩张；取定 \(\mathrm{F}'\) 在 F 上的一组基，就可把 \(\mathrm{F}^{\prime *}\) 与 \(\mathrm{G} = \mathrm{Aut}_{\mathrm{F}}(\mathrm{F}')\) 的一个子群等同起来。设 \(\varphi \in \mathrm{Hom}(\mathrm{F}^{\prime *},\mathrm{K}^{*})\)；计算 \(\mathrm{Ind}_{\mathrm{F}^{\prime *}}^{\mathrm{G}}(\varphi)\)，并证明它等于 \(\mathrm{Ind}_{\mathrm{F}^{\prime *}}^{\mathrm{G}}(\varphi^q)\)。用 \(\lambda\) 表示 \(\varphi\) 在 \(\mathrm{F}^*\) 上的限制，并令 \(\chi_{\varphi} = \chi_{\pi_{\mathrm{P}}\otimes \pi_{\lambda ,1}} - \chi_{\pi_{\lambda ,1}} - \mathrm{Ind}_{\mathrm{F}^{\prime *}}^{\mathrm{G}}(\varphi)\)。证明：若 \(\varphi \neq \varphi^q\)，则我们有 \(\langle \chi_{\varphi},\chi_{\varphi}\rangle = 1\) 与 \(\chi_{\varphi}(1) = q - 1\)，从而 \(\chi_{\varphi}\) 是一个不可约表示 \(\pi_{\varphi}\)（维数为 \(q - 1\) 的表示）的特征标。证明用这种方法我们得到 \(\frac{1}{2} q(q - 1)\) 个两两不同构的不可约表示。

\begin{enumerate}
\def\labelenumi{\alph{enumi})}
\setcounter{enumi}{4}
\tightlist
\item
  证明诸不可约表示 \(\pi_{\lambda}\)（\(\lambda\in\mathrm{Hom}(\mathrm{F}^{*},\mathrm{K}^{*})\)）、\(\pi_{\lambda,\mu}\)（\(\{\lambda,\mu\}\) 是 \(\mathrm{Hom}(\mathrm{F}^{*},\mathrm{K}^{*})\) 的一个二元素子集）与 \(\pi_{\varphi}\)（\(\varphi\in\mathrm{Hom}(\mathrm{F}^{\prime*},\mathrm{K}^{*})\) 且 \(\varphi\neq\varphi^{q}\)）两两不同构，并且 G 的每个不可约表示都同构于它们之一。
\end{enumerate}

\begin{enumerate}
\def\labelenumi{\arabic{enumi})}
\setcounter{enumi}{26}
\tightlist
\item
  设 F 是一个特征为 p \textgreater{} 0 的代数闭体，G 是一个有限群。
\end{enumerate}

\begin{enumerate}
\def\labelenumi{\alph{enumi})}
\item
  设 \(x \in G\)。证明存在唯一的分解 \(x = x_{u}x_{s}\)，其中 \(x_{u}\) 的阶是 p 的幂，\(x_{s}\) 的阶与 p 互素，且 \(x_{u}\) 与 \(x_{s}\) 交换。元素 \(x_{u}\) 与 \(x_{s}\) 都是 x 的幂。
\item
  设 \(\pi\) 是 \(G\) 在 \(F\) 上的一个有限次数线性表示。证明我们有 \(\chi_{\pi}(x) = \chi_{\pi}(x_s)\)（对每个 \(x \in G\)）。
\item
  若 G 的元素 x 的阶与 p 互素，则称它是 p-正则的；用 \(G^{reg}\) 表示 G 的 p-正则元素所成的集合，用 \(\mathcal{Z}_{\mathrm{F}}(\mathrm{G}^{\mathrm{reg}})\) 表示从 \(G^{reg}\) 到 F 的中心函数所成的空间。由同态 \(\mathrm{R}_{\mathrm{F}}(\mathrm{G}) \to \mathcal{Z}_{\mathrm{F}}(\mathrm{G}^{\mathrm{reg}})\)（把表示 \(\pi\) 的类映为 \(\chi_{\pi}\) 在 \(G^{reg}\) 上的限制的同态），我们导出一个环同态 \(\psi : F \otimes_{\mathbf{Z}} R_{\mathrm{F}}(G) \to \mathcal{Z}_{\mathrm{F}}(\mathrm{G}^{\mathrm{reg}})\)。证明 \(\psi\) 是单射（利用 b) 以及 VIII, 第 384 页的推论）。
\item
  设 x 是 G 的一个 p-正则元素，\(Z(x)\) 是它的交换子，P 是 \(Z(x)\) 的一个西罗 p-子群，\(H_{x}\) 是由 P 与 x 生成的 G 的子群。证明 \(H_{x}\) 是 P 与由 x 生成的 G 的子群的直积，并且它的共轭类由 x 的共轭类完全确定。
\item
  构造次数为 1 的表示 \(\sigma_{1},\ldots,\sigma_{m}\)（\(H_{x}\) 的表示）与 F 的元素 \(a_{1},\ldots,a_{m}\)，使得函数 \(\sum a_{i}\chi_{\sigma_{i}}\) 在 x 处取值 1，在 x 的其他幂处取值 0（先构造由 x 生成的子群的表示）。令 \(\rho_{i}=\operatorname{Ind}_{H_{x}}^{G}\pi_{i}\)（对 \(i=1,\ldots,m\)），并令 \(f=\sum_{i}a_{i}\chi_{\rho_{i}}\)。证明 \(f(x)\) 等于 \((\mathrm{Z}(x):\mathrm{H}_{x})1_{\mathrm{F}}\)——它在 F 中非零，并且 \(f(y)\) 对 G 的每个不与 x 共轭的 p-正则元素 y 都为零。
\end{enumerate}

\(f\) ) 由 \(e\) 推出，同态 \(\psi : \mathrm{F} \otimes_{\mathbf{Z}} \mathrm{R}_{\mathrm{F}}(\mathrm{G}) \to \mathcal{Z}_{\mathrm{F}}(\mathrm{G}^{\mathrm{reg}})\) 是双射。特别地，单 F{[}G{]}-模的同构类的个数等于 G 的 \(p\)-正则元素的共轭类的个数。

\begin{enumerate}
\def\labelenumi{\arabic{enumi})}
\setcounter{enumi}{27}
\tightlist
\item
  沿用习题 27 的假设。设 \(m\) 是 \(p\)-正则元素（\(G\) 的元素）的阶的最小公倍数。设 \(F_0\) 是 \(F\) 的由 \(m\) 次单位根生成的子域；它是一个有限域。
\end{enumerate}

\begin{enumerate}
\def\labelenumi{\alph{enumi})}
\item
  证明 G 在 F 上的任何有限次数表示的特征标的所有值都属于 \(F_{0}\)。
\item
  证明 \(F_{0}[G]\) 关于其根的商环是 \(F_{0}\) 上矩阵代数的一个乘积（利用 a) 与韦德伯恩定理）。由此推出，典范同态 \(\mathrm{R}_{\mathrm{F}_{0}}(\mathrm{G}) \to \mathrm{R}_{\mathrm{F}}(\mathrm{G})\) 是双射。
\end{enumerate}

\begin{enumerate}
\def\labelenumi{\arabic{enumi})}
\setcounter{enumi}{28}
\tightlist
\item
  沿用习题 27 与 28 的假设和记号。取定一个本原 \(m\) 次单位根 \(\zeta\)。设 \(\mathcal{O}_m\) 是环 \(\mathbf{Z}[\mathrm{X}] / (\Phi_m)\)（VIII, 第 415 页）；X 在 \(\mathcal{O}_m\) 中的类记作 \(x\)。设 \(\varphi : \mathcal{O}_m \to \mathrm{F}\) 是把 \(x\) 映为 \(\zeta\) 的同态。它诱导出从 \(\mathcal{O}_m\) 中由 \(x\) 生成的子群到群 \(\mu_m(\mathrm{F})\) 的一个同构；用 \(\tau\) 表示逆同构。
\end{enumerate}

\begin{enumerate}
\def\labelenumi{\alph{enumi})}
\item
  设 \(\pi\) 是 G 在 F 上的一个线性表示，维数 \(d\) 有限。对 \(g \in \mathrm{G}\)，设 \(\mathrm{P}_g(\mathrm{X}) = \prod_{i=1}^d (\mathrm{X} - \lambda_i)\) 是 \(\pi(g)\) 的特征多项式；用 \(\beta_{\pi}(g)\) 表示元素 \(\sum_{i} \tau(\lambda_i)\)（\(\mathcal{O}_m\) 的元素）。这样得到的函数 \(\beta_{\pi}: \mathrm{G} \to \mathcal{O}_m\) 称为 \(\pi\) 的布饶尔特征标；我们有 \(\varphi \circ \beta_{\pi} = \chi_{\pi}\) 与 \(\beta_{\pi}(x) = \beta_{\pi}(x_s)\)（对每个 \(x\)），从而 \(\beta_{\pi}\) 由它在 \(\mathrm{G}^{\mathrm{reg}}\) 上的限制确定。
\item
  设 \(\widehat{\mathbf{G}}\) 是单 F{[}G{]}-模的同构类所成的集合。证明诸函数 \((\beta_{\lambda})_{\lambda \in \widehat{\mathbf{G}}}\) 在 \(\mathbf{Z}\) 上线性无关。（若存在一个非平凡关系 \(\sum n_{\lambda} \beta_{\lambda} = 0\)，则可以假设某个 \(n_{\lambda}\) 与 \(p\) 互素，对它应用 \(\varphi\) 并利用习题 27, f) 得到矛盾。）
\item
  设 \(\mathcal{B}\) 是从 \(\mathrm{G}^{\mathrm{reg}}\) 到 \(\mathcal{O}_m\) 的函数所成的集合；同态 \(\mathrm{R}_{\mathrm{F}}(\mathrm{G}) \to \mathcal{B}\)（把 \(\pi\) 的类映为 \(\beta_{\pi}\) 的同态）是单射。由此推出，两个半单 F{[}G{]}-模同构，当且仅当它们的布饶尔特征标相等。
\end{enumerate}

\begin{enumerate}
\def\labelenumi{\arabic{enumi})}
\setcounter{enumi}{29}
\tightlist
\item
  设 \(\mathrm{F}\) 是一个特征为 \(p > 0\) 的代数闭体，\(\mathbf{F}_p\) 是 \(\mathrm{F}\) 的素子域，并设 \(\mathrm{G} = \mathbf{SL}_2(\mathbf{F}_p)\)。用 \(\mathrm{V}\) 表示 F-向量空间 \(\mathrm{F} \otimes_{\mathbf{F}_p} \mathbf{F}_p^2\)，它带有 \(\mathrm{G}\) 的一个作用，该作用由 \(\mathbf{F}_p^2\) 上的作用导出。
\end{enumerate}

\begin{enumerate}
\def\labelenumi{\alph{enumi})}
\item
  证明 F{[}G{]}-模 \(\mathbf{S}^{d}(V)\) 是单模（对 \(0 \leqslant d < p\)）。（设 W 是 F{[}G{]}-模 \(\mathbf{S}^{d}(V)\) 的一个非零子模，\((e, f)\) 是 V 的一组基。证明 \(e^{d}\) 属于 W（利用矩阵 \(\begin{pmatrix}1 & 1 \\ 0 & 1\end{pmatrix}\)），然后证明 W 等于 \(\mathbf{S}^{d}(V)\)（利用矩阵 \(\begin{pmatrix}1 & 0 \\ 1 & 1\end{pmatrix}\)）。）b) 证明 F{[}G{]}-模 \(\mathbf{S}^{p}(V)\) 不是单模（考虑 V 在映射 \(v \mapsto v^{p}\) 下的像）。
\setcounter{enumi}{2}
\item
  对 \(p \geqslant 7\)，证明 \(\mathbf{S}^{p-3}(\mathrm{V})\) 的维数不整除 G 的阶。d) 证明每个单 F{[}G{]}-模都同构于诸模 \(\mathbf{S}^{d}(\mathrm{V})\)（\(0 \leqslant d < p\)）之一（利用习题 27, f)）。
\end{enumerate}

\appendix
\renewcommand{\thechapter}{\arabic{chapter}}
